\section{Methodology}
\label{sec:method}

\subsection{Problem Formulation}

Let $\mathcal{D}=\{x_i\}_{i=1}^{N}$ denote an unlabeled support set sampled from the text distribution. A teacher encoder $T$ maps each input to a cached teacher embedding
\begin{equation}
    t_i = T(x_i) \in \mathbb{R}^{d_T},
\end{equation}
and a compact student encoder $S_{\theta}$ maps the same input to
\begin{equation}
    s_i = S_{\theta}(x_i) \in \mathbb{R}^{d_S}.
\end{equation}
The objective is to train $S_{\theta}$ so that the student preserves the semantic organization induced by the teacher over this finite support, including nearest-neighbor rankings, local semantic regions, and higher-order community structure.

Standard pointwise embedding distillation minimizes a distance between each projected student embedding and its teacher embedding:
\begin{equation}
    \mathcal{L}_{\mathrm{point}}
    =
    \frac{1}{N}
    \sum_{i=1}^{N}
    \left(1-\cos(Ws_i,t_i)\right),
\end{equation}
where $W$ is an optional projection. Relational distillation instead matches student and teacher similarity matrices inside a mini-batch $B$:
\begin{equation}
    \mathcal{L}_{\mathrm{rel}}
    =
    \left\|
    G_B^S - G_B^T
    \right\|_F^2,
\end{equation}
where $G_B^S$ and $G_B^T$ are cosine-similarity matrices over examples in $B$. These objectives are useful local surrogates, but they do not directly supervise the multi-scale neighborhood geometry induced by the teacher.

\subsection{Motivation}

MDD starts from the observation that representation transfer depends on how probability mass is organized over a teacher-induced support graph, not only on the absolute coordinate of each embedding. A compact student may be unable to reproduce the teacher's full embedding map, yet it can still preserve useful semantic geometry if it matches the teacher's neighborhood distributions and global topology.

We therefore treat the teacher embedding space as a semantic graph and distill the random-walk behavior of this graph. For an anchor $x_i$, a one-hop walk captures direct semantic neighbors; multi-hop walks reveal broader regions connected through intermediate examples. This makes the distillation target a soft geometric object: how semantic mass should spread from an anchor to nearby candidates at different scales.

For a diffusion scale $r$, write the full teacher diffusion row as
\begin{equation}
    q_{i,r}(j) = (P_r^T)_{ij}, \qquad j\in\mathcal{D}.
\end{equation}
Since matching a distribution over all $N$ examples is expensive, MDD samples a candidate set $C_i$ and uses the conditional diffusion distribution on that set. Its captured mass is
\begin{equation}
    \rho_{i,r}(C_i)=\sum_{j\in C_i}q_{i,r}(j).
\end{equation}
The method is designed so that $C_i$ contains high-diffusion neighbors, hard negatives, and random negatives, making the restricted target informative while keeping training tractable.

The rest of the section defines the two geometric transfer losses and the stabilizing anchor objective.

\subsection{Multi-Scale Diffusion Loss}

MDD first encodes the support set with the teacher and constructs a teacher-induced support graph. For each example $x_i$, let
\begin{equation}
    \mathcal{N}_k^T(i)
    =
    \operatorname{TopK}_{j\neq i}\cos(t_i,t_j)
\end{equation}
be its top-$k$ teacher nearest neighbors. To reduce noisy one-directional edges and hubness, MDD uses a mutual $k$NN graph:
\begin{equation}
    j \in \mathcal{M}_k^T(i)
    \Longleftrightarrow
    j \in \mathcal{N}_k^T(i)
    \;\mathrm{and}\;
    i \in \mathcal{N}_k^T(j).
\end{equation}
The weighted adjacency matrix is
\begin{equation}
    W_{ij}^{T}
    =
    \begin{cases}
    \exp(\cos(t_i,t_j)/\tau_g), & j \in \mathcal{M}_k^T(i),\\
    0, & \mathrm{otherwise},
    \end{cases}
\end{equation}
where $\tau_g$ controls graph sharpness. Mutual $k$NN can produce isolated support points, so we add self-loops and use a lazy random walk. Let $\bar W^T=W^T+I$ and let $\bar D^T$ be its diagonal degree matrix. The transition matrix is
\begin{equation}
    P^T =
    \frac{1}{2}I
    +
    \frac{1}{2}(\bar D^T)^{-1}\bar W^T.
\end{equation}
This ensures every support point has non-zero degree and stabilizes multi-step diffusion. MDD then defines multi-scale diffusion targets
\begin{equation}
    P_r^T = (P^T)^r,
    \qquad r \in \mathcal{R},
\end{equation}
where $\mathcal{R}$ is typically $\{1,2,4\}$. Small $r$ preserves sharp local neighborhoods, while larger $r$ captures semantic regions and multi-hop communities.

For each anchor $x_i$, MDD builds a single candidate set shared across diffusion scales:
\begin{equation}
    C_i = C_i^{+} \cup C_i^{h} \cup C_i^{r},
\end{equation}
where $C_i^{+}$ contains high-probability diffusion neighbors, $C_i^{h}$ contains teacher hard negatives, and $C_i^{r}$ contains random negatives. The positives are selected from the weighted multi-scale teacher diffusion score
\begin{equation}
    \tilde q_i(j)
    =
    \frac{\sum_{r\in\mathcal{R}}\omega_r q_{i,r}(j)}
    {\sum_{r\in\mathcal{R}}\omega_r},
\end{equation}
as
\begin{equation}
    C_i^{+}
    =
    \operatorname{TopM}_{j\neq i}\tilde q_i(j).
\end{equation}
Hard negatives are teacher-near examples outside the positive set,
\begin{equation}
    C_i^{h}
    =
    \operatorname{TopH}_{j\notin C_i^{+}\cup\{i\}}
    \cos(t_i,t_j),
\end{equation}
and random negatives are sampled uniformly from the remaining support points,
\begin{equation}
    C_i^{r}
    \sim
    \operatorname{Uniform}
    \left(
    \mathcal{D}\setminus(C_i^{+}\cup C_i^{h}\cup\{i\})
    \right).
\end{equation}
Thus the same candidate set contains multi-scale diffusion positives, teacher-similar non-positive examples, and background support points. The teacher diffusion target at scale $r$ is the diffusion mass restricted and renormalized over $C_i$:
\begin{equation}
    p_{i,r}^{T}(j)
    =
    \frac{(P_r^T)_{ij}}
    {\sum_{j' \in C_i}(P_r^T)_{ij'}},
    \qquad j\in C_i.
\end{equation}

Given an anchor embedding $s_i$ and candidate embeddings $\{s_j:j\in C_i\}$, the student induces a distribution over the same candidate set:
\begin{equation}
    p_i^{S}(j)
    =
    \frac{\exp(\cos(s_i,s_j)/\tau_s)}
    {\sum_{j'\in C_i}\exp(\cos(s_i,s_{j'})/\tau_s)},
    \qquad j\in C_i,
\end{equation}
where $\tau_s$ is the student temperature. The diffusion loss is
\begin{equation}
    \mathcal{L}_{\mathrm{diff}}
    =
    \frac{1}{|B|}
    \sum_{i\in B}
    \sum_{r\in\mathcal{R}_{e}}
    \omega_r
    D_{\mathrm{KL}}
    \left(
        p_{i,r}^{T} \| p_i^{S}
    \right),
\end{equation}
where $\omega_r$ is the scale weight and $\mathcal{R}_{e}$ is the set of active scales at epoch $e$.

Although the teacher target varies with $r$, the student distribution $p_i^S$ is intentionally shared across active scales. The weighted objective is equivalent to matching a multi-scale diffusion barycenter. Define
\begin{equation}
    \tilde p_i^T(j)
    =
    \frac{\sum_{r\in\mathcal{R}_e}\omega_r p_{i,r}^T(j)}
    {\sum_{r\in\mathcal{R}_e}\omega_r},
    \qquad j\in C_i.
\end{equation}
Then
\begin{equation}
    \sum_{r\in\mathcal{R}_e}
    \omega_r
    D_{\mathrm{KL}}
    \left(p_{i,r}^{T}\|p_i^{S}\right)
    =
    \Omega_e
    D_{\mathrm{KL}}
    \left(\tilde p_i^T\|p_i^{S}\right)
    +
    \mathrm{const.},
\end{equation}
where $\Omega_e=\sum_{r\in\mathcal{R}_e}\omega_r$ and the constant does not depend on the student. Thus MDD trains the student to match one weighted multi-scale diffusion profile rather than forcing it to satisfy incompatible scale-specific distributions.

\subsection{Spectral Manifold Anchoring}

Diffusion matching transfers local and multi-hop neighborhoods, but a student can still fit candidate distributions while distorting broader community structure. MDD therefore computes spectral coordinates from the teacher graph. Let the symmetric normalized graph Laplacian be
\begin{equation}
    L^T = I - (\bar D^T)^{-1/2}\bar W^T(\bar D^T)^{-1/2}.
\end{equation}
Let $U^T\in\mathbb{R}^{N\times m}$ contain the first $m$ non-trivial eigenvectors of $L^T$, and let $u_i^T$ be the row corresponding to $x_i$. A learnable projection maps student embeddings to this spectral space:
\begin{equation}
    z_i^S = W_{\mathrm{spec}}s_i.
\end{equation}
The spectral loss is
\begin{equation}
    \mathcal{L}_{\mathrm{spec}}
    =
    \frac{1}{|B|}
    \sum_{i\in B}
    \left\|
    z_i^S-u_i^T
    \right\|_2^2.
\end{equation}

The first non-trivial eigenvectors of the normalized graph Laplacian are low-frequency modes of the teacher semantic graph. They vary smoothly over strongly connected regions and separate broader communities, following the graph-manifold principles behind Laplacian Eigenmaps and Diffusion Maps~\citep{belkin2003laplacian,coifman2006diffusionmaps}. The spectral loss therefore anchors the student to the teacher's global topology.

\subsection{Optimization Objectives}

MDD includes a lightweight pointwise anchor for stability:
\begin{equation}
    \mathcal{L}_{\mathrm{anchor}}
    =
    \frac{1}{|B|}
    \sum_{i\in B}
    \left(1-\cos(W_a s_i,t_i)\right),
\end{equation}
where $W_a$ maps student embeddings to the teacher dimension. This term acts as a boundary condition on the embedding map, while the diffusion and spectral objectives carry the main geometric supervision.

The final objective is
\begin{equation}
\begin{aligned}
    \mathcal{L}_{\mathrm{MDD}}
    ={}&
    \lambda_{1}\mathcal{L}_{\mathrm{diff}}
    +
    \lambda_{2}\mathcal{L}_{\mathrm{spec}}
    \\
    &+
    \lambda_{3}\mathcal{L}_{\mathrm{anchor}}.
\end{aligned}
\end{equation}
We use a diffusion-scale curriculum: the first epoch uses one-step diffusion, the second epoch adds two-step diffusion, and later epochs add the full scale set and spectral anchoring. This avoids forcing long-range, smoother structure before the student has learned reliable local neighborhoods.

\subsection{Theoretical Analysis}

We provide a finite-support analysis showing that MDD is a controlled surrogate for preserving teacher-induced diffusion geometry. The analysis separates four effects: multi-scale aggregation, candidate truncation, low-frequency spectral anchoring, and the coverage gap of purely mini-batch relational matching.

\paragraph{Multi-scale diffusion as barycenter matching.}
For an anchor $x_i$, let $\alpha_r=\omega_r/\sum_{r'\in\mathcal{R}_e}\omega_{r'}$ be the normalized active scale weights and define the multi-scale diffusion barycenter
\begin{equation}
    \bar p_i^T(j)=
    \sum_{r\in\mathcal{R}_e}\alpha_r p_{i,r}^T(j),
    \qquad j\in C_i.
\end{equation}
\noindent\textbf{Proposition 1.}
For any student distribution $p_i^S$ over $C_i$,
\begin{equation}
\begin{aligned}
&\sum_{r\in\mathcal{R}_e}\alpha_r
D_{\mathrm{KL}}(p_{i,r}^T\|p_i^S)\\
&\quad=
D_{\mathrm{KL}}(\bar p_i^T\|p_i^S)
\\
&\qquad+
\sum_{r\in\mathcal{R}_e}\alpha_r
D_{\mathrm{KL}}(p_{i,r}^T\|\bar p_i^T).
\end{aligned}
\end{equation}
The second term is independent of the student. Thus, optimizing the multi-scale KL loss is equivalent to matching the student to the teacher's multi-scale diffusion barycenter, up to a constant measuring scale diversity among teacher neighborhoods.

\emph{Proof sketch.}
Expand each KL term as
\begin{equation}
    \sum_j p_{i,r}^T(j)\log p_{i,r}^T(j)
    -
    \sum_j p_{i,r}^T(j)\log p_i^S(j).
\end{equation}
Adding and subtracting the barycenter log-density gives the stated decomposition.

\paragraph{Candidate-truncated diffusion matching.}
Let $q_{i,r}$ be the full teacher diffusion row over $\mathcal{D}$, let
\begin{equation}
    \rho_{i,r}(C_i)=\sum_{j\in C_i}q_{i,r}(j),
\end{equation}
and let $q_{i,r}^{C}$ be the conditional distribution of $q_{i,r}$ on $C_i$. Let $\widehat p_i^S$ be the zero extension of $p_i^S$ from $C_i$ to $\mathcal{D}$.

\noindent\textbf{Theorem 1.}
Assume $\rho_{i,r}(C_i)\ge 1-\epsilon$ for all $r\in\mathcal{R}_e$. If
\begin{equation}
    \sum_{r\in\mathcal{R}_e}\alpha_r
    D_{\mathrm{KL}}(q_{i,r}^{C}\|p_i^S)
    \le \eta,
\end{equation}
then
\begin{equation}
    \sum_{r\in\mathcal{R}_e}\alpha_r
    \left\|q_{i,r}-\widehat p_i^S\right\|_1
    \le
    2\epsilon+\sqrt{2\eta}.
\end{equation}
Consequently, for any bounded probe function $f:\mathcal{D}\rightarrow[-1,1]$,
\begin{equation}
\begin{aligned}
&\sum_{r\in\mathcal{R}_e}\alpha_r
\left|
\mathbb{E}_{j\sim q_{i,r}}[f(j)]
-
\mathbb{E}_{j\sim p_i^S}[f(j)]
\right|\\
&\quad\le
2\epsilon+\sqrt{2\eta}.
\end{aligned}
\end{equation}
This separates the approximation error into candidate truncation error $\epsilon$ and student matching error $\eta$.

\emph{Proof sketch.}
For each scale,
\begin{equation}
    \left\|q_{i,r}-\widehat p_i^S\right\|_1
    \le
    2(1-\rho_{i,r}(C_i))
    +
    \left\|q_{i,r}^{C}-p_i^S\right\|_1.
\end{equation}
Pinsker's inequality bounds the second term by
\begin{equation}
    \sqrt{2D_{\mathrm{KL}}(q_{i,r}^{C}\|p_i^S)}.
\end{equation}
Averaging over scales and applying Jensen's inequality gives the result. The probe bound follows from $|f|\le 1$.

\paragraph{Low-frequency topology preservation.}
Let $u_i^T$ be the teacher spectral coordinate and $z_i^S=W_{\mathrm{spec}}s_i$ be the student spectral prediction. Suppose the support-average spectral loss satisfies
\begin{equation}
    \frac{1}{N}\sum_{i=1}^{N}
    \left\|z_i^S-u_i^T\right\|_2^2
    \le \xi.
\end{equation}
Define $d_{T,m}(i,j)=\|u_i^T-u_j^T\|_2$ and
$d_{S,m}(i,j)=\|z_i^S-z_j^S\|_2$.

\noindent\textbf{Theorem 2.}
If $(i,j)$ is sampled uniformly from $\mathcal{D}\times\mathcal{D}$, then
\begin{equation}
    \mathbb{E}_{i,j}
    \left[
    \left|d_{S,m}(i,j)-d_{T,m}(i,j)\right|
    \right]
    \le
    2\sqrt{\xi}.
\end{equation}
More generally, if the spectral coordinates are multiplied by diagonal diffusion weights whose largest retained weight is $a_{\max}$, the bound becomes $2a_{\max}\sqrt{\xi}$.

\emph{Proof sketch.}
The reverse triangle inequality gives
\begin{equation}
\begin{aligned}
&\left|d_{S,m}(i,j)-d_{T,m}(i,j)\right|\\
&\quad\le
\left\|z_i^S-u_i^T\right\|_2
+
\left\|z_j^S-u_j^T\right\|_2.
\end{aligned}
\end{equation}
Taking expectation and applying Jensen's inequality gives the result. This guarantee concerns preservation of teacher low-frequency coordinates encoded by the student; it does not claim that the student-induced graph recovers the full teacher spectrum.

\paragraph{Mini-batch relational coverage gap.}
Let $B$ be a batch of size $b$ containing anchor $i$, with the other $b-1$ examples sampled uniformly without replacement from $\mathcal{D}\setminus\{i\}$. For any teacher diffusion row $q_{i,r}$,
\begin{equation}
    \mathbb{E}_{B}
    \left[
    \sum_{j\in B\setminus\{i\}}q_{i,r}(j)
    \right]
    =
    \frac{b-1}{N-1}
    \left(1-q_{i,r}(i)\right).
\end{equation}
When $b\ll N$, a batch-local relational objective observes only a small expected fraction of the diffusion neighborhood around each anchor. MDD avoids this coverage bottleneck by constructing $C_i$ from high-diffusion neighbors, teacher hard negatives, and random negatives.

\begin{algorithm}[t]
\caption{MDD Training}
\label{alg:mddkd}
\begin{algorithmic}[1]
\REQUIRE Support set $\mathcal{D}$, teacher $T$, student $S_{\theta}$, graph size $k$, diffusion scales $\mathcal{R}$
\STATE Cache teacher embeddings $t_i=T(x_i)$ for all $x_i\in\mathcal{D}$
\STATE Build a mutual $k$NN teacher graph $W^T$
\STATE Add self-loops and compute lazy transition matrix $P^T$
\STATE Precompute sparse diffusion candidates and targets $p_{i,r}^T$
\STATE Compute spectral coordinates $u_i^T$ from the graph Laplacian
\FOR{each training epoch}
    \STATE Activate diffusion scales according to the curriculum
    \FOR{each mini-batch $B$}
        \STATE Encode anchors and candidate texts with $S_{\theta}$
        \STATE Compute student candidate distributions $p_i^S$
        \STATE Compute $\mathcal{L}_{\mathrm{diff}}$, $\mathcal{L}_{\mathrm{spec}}$, and $\mathcal{L}_{\mathrm{anchor}}$
        \STATE Update $\theta$ by minimizing $\mathcal{L}_{\mathrm{MDD}}$
    \ENDFOR
\ENDFOR
\RETURN Distilled student encoder $S_{\theta}$
\end{algorithmic}
\end{algorithm}
